\clearpage
\section{第 7.1 节：为什么研究谐振子}
\label{section:QM-C07-S01}

\studymeta
  {QM-C07-S01}
  {2026-09-19}
  {Shankar, \textit{Principles of Quantum Mechanics}, Second Edition}
  {第 7 章 \S 7.1；印刷页 pp.~185--188；PDF pp.~197--200}
  {讲解与概念核对已完成；笔记已确认}

\textbf{范围与主线。}本节解释谐振子为何反复出现在物理问题中：稳定平衡附近的非简并小振动近似为单个谐振子；多个耦合自由度可在二次近似下分解为独立的正常模；晶体振动和自由电磁场提供了多模系统的例子。本节回顾第 1.8 节的教材例 1.8.6，其完整题干和解答已收于\relatedexercise{QM-C01-S08-E06}{课程题目部分}，此处不重复。四次势及零、负曲率的边界属于严格性补充。

\notetopic{QM-C07-S01-T01}{稳定平衡附近的二次近似}

\keyidea{谐振子不是任意势能的精确模型；它来自平衡点附近保留最低非零二次项的局部近似。}

令 $x_0$ 为光滑势能 $V(x)$ 的平衡位置，$q=x-x_0$ 为小位移。若 $x_0$ 是内点，则 $V'(x_0)=0$。在足够光滑的条件下，
\begin{align}
  V(x_0+q)
  &=V(x_0)+V'(x_0)q+\frac12V''(x_0)q^2
    +\frac16V'''(x_0)q^3+\cdots\\
  &=V(x_0)+\frac12k_{\mathrm{eff}}q^2+O(q^3),
  \qquad k_{\mathrm{eff}}:=V''(x_0).
  \label{eq:QM-C07-S01-taylor}
\end{align}
势能常数 $V(x_0)$ 不影响经典运动，可精确地调整能量零点；线性项因平衡条件精确消失。只有舍去三次及更高次项是近似，要求在所考察的位移范围内这些项相对二次项足够小。例如三次项非零时，可检查
\begin{equation}
  \left|\frac{V'''(x_0)q}{3V''(x_0)}\right|\ll1,
\end{equation}
同时还须检查更高阶项。

若 $k_{\mathrm{eff}}>0$，该平衡点在二次阶非简并。于是
\begin{equation}
  m\ddot q=-V'(x_0+q)\simeq-k_{\mathrm{eff}}q,
  \qquad
  \boxed{\ddot q+\omega^2q\simeq0},
  \qquad
  \boxed{\omega=\sqrt{\frac{V''(x_0)}{m}}}.
  \label{eq:QM-C07-S01-local-frequency}
\end{equation}
若势能本来就是二次函数，方程精确成立；对一般势能，它只描述平衡附近的小振动。此处是\emph{经典}小振动的动机，尚未给出量子谐振子的能谱。

\textbf{适用边界。}稳定平衡不必满足 $V''(x_0)>0$：例如 $V(q)=\lambda q^4$（$\lambda>0$）在 $q=0$ 稳定，但二次项为零，不能由式~\eqref{eq:QM-C07-S01-local-frequency} 得到非零的谐振频率。若某方向的二次曲率为负，该方向的平衡则不稳定，见下文的矩阵判据。

\notetopic{QM-C07-S01-T02}{两个耦合振子：换到正常坐标}

\keyidea{原坐标中的耦合可以来自坐标选择；对角化二次势能后，每个正常坐标独立振动。}

以教材例 1.8.6 中两个等质量物块、三个相同弹簧为例。设弹簧劲度系数为 $k$，$\omega_0^2=k/m$，则
\begin{equation}
  H=\frac{p_1^2+p_2^2}{2m}
    +\frac{k}{2}\bigl[x_1^2+(x_1-x_2)^2+x_2^2\bigr]
   =\frac{\mathbf p^{\mathsf T}\mathbf p}{2m}
    +\frac12\mathbf x^{\mathsf T}K\mathbf x,
  \qquad
  K=k\begin{pmatrix}2&-1\\-1&2\end{pmatrix}.
\end{equation}
对称矩阵 $K$ 的两个正交归一特征向量是 $(1,1)^{\mathsf T}/\sqrt2$ 和 $(1,-1)^{\mathsf T}/\sqrt2$，特征值分别为 $k$ 和 $3k$。据此取
\begin{equation}
  Q_I=\frac{x_1+x_2}{\sqrt2},\qquad
  Q_{II}=\frac{x_1-x_2}{\sqrt2},\qquad
  P_I=\frac{p_1+p_2}{\sqrt2},\qquad
  P_{II}=\frac{p_1-p_2}{\sqrt2}.
\end{equation}
该变换正交，所以 $p_1^2+p_2^2=P_I^2+P_{II}^2$；将 $x_1=(Q_I+Q_{II})/\sqrt2$、$x_2=(Q_I-Q_{II})/\sqrt2$ 代入势能，得到
\begin{equation}
  x_1^2+(x_1-x_2)^2+x_2^2=Q_I^2+3Q_{II}^2.
\end{equation}
于是
\begin{equation}
  \boxed{H=\left(\frac{P_I^2}{2m}+\frac{kQ_I^2}{2}\right)
       +\left(\frac{P_{II}^2}{2m}+\frac{3kQ_{II}^2}{2}\right)},
  \qquad
  \boxed{\omega_I=\omega_0,\quad\omega_{II}=\sqrt3\,\omega_0}.
  \label{eq:QM-C07-S01-two-modes}
\end{equation}
当 $Q_{II}=0$ 时 $x_1=x_2$，中间弹簧不伸缩；只激发 $Q_{II}$ 时 $x_1=-x_2$，中间弹簧参与恢复力。$3\omega_0^2$ 是频率的\emph{平方}，因此频率为 $\sqrt3\,\omega_0$。模振幅 $Q_{II}$ 本身可正可负，其符号与特征值 $3k>0$ 无关。教材例题中从运动方程求正常模及位移传播算符的推导见\relatedexercise{QM-C01-S08-E06}{例 1.8.6}。

\notetopic{QM-C07-S01-T03}{多自由度系统与正常模的适用范围}

\keyidea{在二次近似下，一个耦合系统可以改写为独立振子的集合；每个模是整个系统的一种集体运动图样。}

对于 $N$ 个等质量位移自由度，令 $\mathbf q$ 表示相对平衡位置的位移。平衡点处的一阶导数为零，势能的二次近似为
\begin{equation}
  V(\mathbf q)\simeq V(\mathbf0)+\frac12\mathbf q^{\mathsf T}K\mathbf q,
  \qquad
  K_{ij}=\left.\frac{\partial^2V}{\partial q_i\partial q_j}\right|_{\mathbf q=\mathbf0}.
\end{equation}
在所需光滑性下，$K$ 是实对称矩阵。取正交矩阵 $O$，使
\begin{equation}
  O^{\mathsf T}KO=\operatorname{diag}(\kappa_1,\ldots,\kappa_N),
  \qquad \mathbf Q=O^{\mathsf T}\mathbf q,\quad
  \mathbf P=O^{\mathsf T}\mathbf p.
\end{equation}
舍去势能常数并保留二次项后，
\begin{equation}
  H\simeq\sum_{\alpha=1}^{N}
     \left(\frac{P_\alpha^2}{2m}+\frac12\kappa_\alpha Q_\alpha^2\right),
  \qquad
  \ddot Q_\alpha+\frac{\kappa_\alpha}{m}Q_\alpha\simeq0.
  \label{eq:QM-C07-S01-general-modes}
\end{equation}
若 $\kappa_\alpha>0$，该模稳定振动，频率为 $\sqrt{\kappa_\alpha/m}$。若 $\kappa_\alpha=0$，二次近似在该方向没有恢复力，是否真正自由还取决于更高阶项或对称性。若 $\kappa_\alpha<0$，则 $\ddot Q_\alpha=(|\kappa_\alpha|/m)Q_\alpha$，小偏离一般呈指数增长，不能解释为稳定振动；负的是势能曲率，而非坐标 $Q_\alpha$ 的取值。

\textbf{不同质量及非线性。}不同质量时，动能不是 $\mathbf p^{\mathsf T}\mathbf p/(2m)$，须使用质量矩阵或质量加权坐标再对角化。舍去的三次及更高阶项会耦合不同正常模，因此“独立振子”一般只在小振动的谐近似内成立。

\textbf{晶体与场。}含 $N_0$ 个可在三维空间位移的原子的有限晶体，在谐近似下有 $3N_0$ 个位移自由度，因而可分解为相应数量的正常模。在具有平移对称性的理想晶体中，可进一步用波矢和分支标记集体振动；波矢描述空间振动图样，而不是某个原子的行进轨迹。自由电磁场在选取物理横向自由度后，也可按波矢及两个横向偏振分解为模式。晶体的具体分支结构和场的边界条件需分别处理，不能把简单模型中的“一个纵模、两个横模”当作任意晶体在所有波矢处的普遍结论。

\subsection{一分钟回顾}
平衡附近出现谐振子，需要非零正二次曲率以及足够小的位移。耦合系统的正常模由二次势能矩阵的特征向量给出；特征值决定频率的平方，不决定模振幅的正负。晶体振动与自由场说明了为何要先研究单个振子，再研究许多振子的组合。

\section{第 7.2 节：经典谐振子回顾}
\label{section:QM-C07-S02}

\studymeta
  {QM-C07-S02}
  {2026-09-19 至 2026-09-20}
  {Shankar, \textit{Principles of Quantum Mechanics}, Second Edition}
  {第 7 章 \S 7.2；印刷页 pp.~188--189；PDF pp.~200--201}
  {讲解与概念核对已完成；笔记待确认}

\textbf{范围与主线。}本节复习一维经典谐振子的运动与能量，为下一节比较量子振子作准备。教材此处没有独立编号的 Example 或 Exercise；速度符号的说明是对式 (7.2.5) 的严格性补充。

\notetopic{QM-C07-S02-T01}{运动方程、初始条件与连续能量}

\keyidea{经典谐振子的能量由振幅决定且连续可取；同一能量和位置通常对应两个相反方向的动量。}

考虑无驱动、无阻尼的振子，取 $m>0$、$\omega>0$，其 Hamiltonian 为
\begin{equation}
  \mathcal H(x,p)=\frac{p^2}{2m}+\frac12m\omega^2x^2.
\end{equation}
Hamilton 方程及由此得到的二阶方程是
\begin{equation}
  \dot x=\pdv{\mathcal H}{p}=\frac{p}{m},\qquad
  \dot p=-\pdv{\mathcal H}{x}=-m\omega^2x,
  \qquad
  \boxed{\ddot x+\omega^2x=0}.
  \label{eq:QM-C07-S02-motion}
\end{equation}
设初始数据为 $x_i=x(0)$、$p_i=p(0)$。通解可以写成
\begin{align}
  x(t)&=A\cos\omega t+B\sin\omega t
       =x_0\cos(\omega t+\phi),\\
  A&=x_i,\qquad B=\frac{p_i}{m\omega},\qquad
  x_0=\sqrt{x_i^2+\left(\frac{p_i}{m\omega}\right)^2}\ge0.
  \label{eq:QM-C07-S02-solution}
\end{align}
这里教材的 $x_0$ 表示\emph{振幅}，一般并非初始位置 $x_i$；它与相位 $\phi$ 一起由初始位置和动量确定。对应动量为
\begin{equation}
  p(t)=m\dot x(t)=-m\omega x_0\sin(\omega t+\phi).
\end{equation}
代入 $\mathcal H$，并利用 $\sin^2+\cos^2=1$，得到守恒的能量
\begin{equation}
  \boxed{E=\frac12m\dot x^2+\frac12m\omega^2x^2
         =\frac12m\omega^2x_0^2}.
  \label{eq:QM-C07-S02-energy}
\end{equation}
因为经典振幅 $x_0$ 可连续取非负值，所以经典能量可连续取 $E\ge0$；最低能量 $E=0$ 对应 $x=p=0$，即始终静止于平衡位置。

\textbf{位置、速率与运动方向。}由能量式解出
\begin{equation}
  \dot x^2=\frac{2E}{m}-\omega^2x^2
            =\omega^2(x_0^2-x^2),
  \qquad
  \boxed{\dot x=\pm\omega\sqrt{x_0^2-x^2}}
  \quad (\abs{x}\le x_0).
  \label{eq:QM-C07-S02-velocity-branches}
\end{equation}
在 $\abs{x}<x_0$ 且 $E>0$ 时，正负号分别表示向右和向左运动；两个转折点 $x=\pm x_0$ 的速度为零，原点处的速率最大。教材式 (7.2.5) 只写正平方根，应将其理解为速率 $\abs{\dot x}$，或理解为 $\dot x\ge0$ 的分支，而非整个周期的有符号速度。

\subsection{一分钟回顾}
Hamilton 方程给出 $\ddot x+\omega^2x=0$。振幅由位置与动量两项初始数据共同决定，能量为 $m\omega^2x_0^2/2$ 且在经典理论中连续。同一能量和内部位置对应两种运动方向；下一节将以这些结论对照量子谐振子。
